\section{Case Studies: Ithaca, Ann Arbor, and Champaign-Urbana}
\label{sec:casestudies}

\subsection{Ithaca, New York: The Cornell Effect}

Tompkins County, New York, home to Cornell University (enrollment $\approx$25,000) and Ithaca College ($\approx$7,000), recorded a 2020 Census population of approximately 105,740. Of this total, an estimated 28,000--30,000 residents are college students counted under the campus rule. The non-student resident population is therefore approximately 75,000--77,000---a figure that would place Tompkins County well below the threshold necessary to anchor a standalone congressional district.

Under New York's post-2010 redistricting geometry, Tompkins County is combined with neighboring counties to form a congressional district. The student population contributes meaningfully to the district's population total, affecting which combination of surrounding counties is required to reach the target population ($\approx$776,000 in the 2022 cycle). A home-address correction removing 28,000 Tompkins County students would require the district to draw more of Chemung, Schuyler, or Tioga counties to meet the equipopulation target---potentially changing the district's partisan composition.

\textbf{Litigation context}: New York's student-counting litigation has primarily addressed state legislative redistricting. In \emph{Finch v. New York} and related proceedings, plaintiffs challenged the assignment of student populations without verification of home-address status. The New York Court of Appeals has generally upheld campus counting for state redistricting, but the legislature's 2010 prison-gerrymandering reform (counting prisoners at home address) was not extended to college students, creating an asymmetry: prisoners are counted at home, students at campus, with opposing partisan effects.

\subsection{Ann Arbor, Michigan: The University of Michigan Effect}

Washtenaw County, Michigan, home to the University of Michigan (enrollment $\approx$47,000), recorded a 2020 Census population of approximately 367,601. The student population represents approximately 14.6\% of the county total---a substantial fraction, though diluted by the large non-student resident population in the broader county.

At the census-tract level, the effect is more concentrated. The central Ann Arbor tracts containing the Central Campus, South Quad, and West Quad dormitory complexes have student populations constituting 60--85\% of Census enumeration. These tracts' populations are allocated to Michigan's 6th Congressional District (2022 cycle), which covers Ann Arbor and portions of Washtenaw County.

Under home-address reallocation, approximately 35,000--40,000 students would be redistributed to their pre-enrollment Michigan home counties---predominantly Oakland, Macomb, Kent, and Ingham counties---and to out-of-state origins. The 6th District would lose 35,000--40,000 persons, requiring it to absorb additional territory from adjacent districts. Given that Ann Arbor is surrounded by more Republican-leaning suburban territory (Livingston County to the east, Lenawee County to the south), this expansion would dilute the district's Democratic lean.

\subsection{Champaign-Urbana, Illinois: The UIUC Effect}

The University of Illinois Urbana-Champaign (enrollment $\approx$56,000) is the most extreme case among major public research universities. Champaign County's 2020 Census population of approximately 209,448 includes an estimated 55,000--58,000 students, constituting approximately 27--28\% of the county total. Combined with the Urbana portions of the twin-city area, the student fraction exceeds one-third of Census population in the core Champaign-Urbana tracts.

Illinois's 13th Congressional District (2022 cycle) is centered on Champaign-Urbana. The district's liberal Democratic lean in recent elections is substantially attributable to student voting and student population credit. Under home-address reallocation, removing 55,000 students from the district's population total would require absorbing substantially more Republican-leaning downstate Illinois territory, almost certainly producing a competitive or Republican-leaning district from what is currently a reliably Democratic seat.

This case illustrates the strongest version of the college-town inflation argument: the district's partisan character---indeed, its very configuration as a viable Democratic-leaning unit---depends critically on the decision to count 55,000 students at campus rather than at their home addresses, the majority of which are in the Chicago suburban counties that send students to UIUC in large numbers.

\subsection{Pattern Across Cases}

The three case studies exhibit a common pattern:
\begin{enumerate}
    \item A major research university with $>$25,000 enrollment is surrounded by a relatively small permanent resident community.
    \item Campus counting adds a large student population that is politically distinct from the surrounding permanent residents.
    \item Students vote (when they vote) at far higher Democratic rates than the surrounding non-student population.
    \item The students' home communities---from which they originated and to which most will return---are suburban counties with a different partisan character.
\end{enumerate}

The directional asymmetry is consistent across all three cases: campus counting benefits the college-town congressional district (Democratic-leaning) and reduces the population available to adjacent suburban districts (Republican-leaning).
